\documentclass[12pt]{article}
\usepackage[letterpaper,margin=1in]{geometry}
\usepackage{fontspec}
\usepackage{setspace}
\usepackage{amsmath}
\usepackage{amssymb}
\usepackage{microtype}
\usepackage{xcolor}
\usepackage{fancyhdr}
\usepackage{titlesec}
\usepackage{parskip}
\usepackage{hyperref}
\hypersetup{
  colorlinks=true,
  linkcolor=blue!50!black,
  urlcolor=blue!50!black,
  citecolor=blue!50!black,
  pdftitle={M2 to A3: The Phantastic Motion and the Phantasma Proper},
  pdfauthor={Dalton Salvo},
  pdfsubject={Dissertation -- Aristotelian Actualization Chain}
}

% Main font: Linux Libertine if available, otherwise default serif
\IfFontExistsTF{Linux Libertine O}{%
  \setmainfont{Linux Libertine O}[Numbers=OldStyle]%
}{%
  \IfFontExistsTF{TeX Gyre Pagella}{%
    \setmainfont{TeX Gyre Pagella}%
  }{}%
}

% Greek wrapping: try dedicated Greek font, else fall back to main font
\IfFontExistsTF{GFS Didot}{%
  \newfontfamily{\greekfont}{GFS Didot}%
  \newcommand{\gk}[1]{{\greekfont #1}}%
}{%
  \IfFontExistsTF{DejaVu Serif}{%
    \newfontfamily{\greekfont}{DejaVu Serif}%
    \newcommand{\gk}[1]{{\greekfont #1}}%
  }{%
    % Last resort: use main font (works if it has Greek glyphs)
    \newcommand{\gk}[1]{#1}%
  }
}

% Line spacing and paragraph style
\setstretch{1.3}
\setlength{\parindent}{0pt}
\setlength{\parskip}{0.8em}

% Section titles: generous spacing, clean typography
\titleformat{\section}{\Large\bfseries}{}{0pt}{}
\titleformat{\subsection}{\large\bfseries}{}{0pt}{}
\titlespacing*{\section}{0pt}{2em}{1em}
\titlespacing*{\subsection}{0pt}{1.5em}{0.6em}

% Footnote style: slightly tighter
\usepackage[hang,flushmargin]{footmisc}

% Page header/footer
\setlength{\headheight}{14pt}
\pagestyle{fancy}
\fancyhf{}
\fancyhead[L]{\small\itshape M\textsubscript{2}$\to$A\textsubscript{3}: Phantasma Proper}
\fancyhead[R]{\small\itshape Draft — \today}
\fancyfoot[C]{\thepage}
\renewcommand{\headrulewidth}{0.4pt}
\renewcommand{\footrulewidth}{0pt}

% Title page content
\title{\vspace{-2em}\textbf{$M_2 \to A_3$: The Phantastic Motion and the Phantasma Proper}\\[0.5em]
       \large\normalfont A Dissertation Section (Console Draft)}
\author{Dalton Salvo}
\date{April 2026 — Draft for Review}

\begin{document}
\maketitle

\thispagestyle{fancy}

\subsection*{1. $M_2$→$A_3$: The Phantastic Motion and the Phantasma Proper}

$A_2$ closed with the perceptual motion complete and a residual motion --- resonant \textit{kinēsis} in its dual-aspect structure of resonant \textit{aisthēma} and resonant \textit{orexis} --- deposited in the sensory apparatus. The present section specifies how that residue becomes the \textit{phantasma}. $M_2$$\to$$A_3$ is the motion through which \textit{phantasia}, as a distinct psychic faculty, takes up the resonant \textit{kinēsis} left by perception and crystallizes it into a determinate image available to the soul as the content of memory, imagination, dreaming, deliberation, and thought. Its terminus $A_3$ --- the \textit{phantasma} proper --- is an actuality ontologically distinct from the residue upon which it depends: it is the \textit{energeia} of a secondary motion whose material cause is the dual trace at $A_2$ but whose formal reality is not reducible to that trace.

The section proceeds through seven architectural moves. First, it establishes the distinction between \textit{phantasia} as \textit{kinēsis} (the ongoing activity) and the \textit{phantasma} as its formal product. Second, it applies the three-factor kinetic schema to the phantastic motion, identifying the residual trace as unmoved originator, \textit{phantasia} as moved mover, and the central perceptual organ --- the \textit{koinē aisthēsis} --- as that which is moved. Third, it analyses the \textit{phantasma} as the crystallized actuality produced at the terminus of this secondary motion. Fourth, it argues that the \textit{phantasma} inherits the dual-trace structure deposited at $A_2$, coming into being already charged with both formal-epistemic content and affective-orectic valence. Fifth, it develops Caston's causal-layering analysis to show that the \textit{phantasma}, though causally downstream of the residue, is ontologically distinct from it --- the \textit{phantasma} can diverge from its causal ancestry in ways the residue cannot. Sixth, it specifies the \textit{phantasma}'s role as the object upon which subsequent cognitive and orectic motions operate, grounding the "the soul never thinks without an image" thesis of \textit{De Anima} III.7-8. Finally, the synthesis situates $A_3$ as the unmoved originator of the cognitive motions at $M_3$$\to$$M_4$ that produce the completed cognitive actualities (intellection, memory, deliberation, and the committal dimension of \textit{doxa}) whose analysis is the subject of the next section.

Two architectural stakes govern the analysis. First, the distinction between the residue (resonant \textit{kinēsis}) and the \textit{phantasma} is not merely terminological but ontological: if the two were identified, \textit{phantasia} would collapse into a mere persistence of perception, and its distinctive cognitive functions --- dreaming, imagining what is absent, recombining elements into what has never been perceived, presenting objects for deliberation --- would become inexplicable. Second, the \textit{phantasma}'s inheritance of the dual-trace structure is what licenses everything downstream: without the affective charge inherited from resonant \textit{orexis}, the \textit{phantasma} could not move \textit{orexis}; without the formal content inherited from resonant \textit{aisthēma}, it could not serve as the object of thought. $A_3$ is the hinge at which the physical-perceptual register of $A_1$$\to$$A_2$ is transformed into the cognitive-affective register that the rest of the chain elaborates.

\subsection*{2. Phantasia as Kinesis, Phantasma as Product}

Aristotle's clearest definition of \textit{phantasia} appears at \textit{De Anima} III.3, 428b10-17: "But since when one thing has been set in motion another thing may be moved by it, and imagination is held to be a movement and to be impossible without sensation, i.e. to occur in beings that are percipient and to have for its content what can be perceived, and since movement may be produced by actual sensation and that movement is necessarily similar in character to the sensation itself, this movement cannot exist apart from sensation or in creatures that do not perceive, and its possessor does and undergoes many things in virtue of it, and it is true and false." Four features are stipulated. \textit{Phantasia} is a motion (\textit{kinēsis}). It is produced by actual sensation and similar in character to the sensation that produced it. It requires a percipient being. And its products admit of truth and falsity --- a feature Aristotle treats as decisive for distinguishing \textit{phantasia} from perception of proper sensibles, which is "always true or admits the least falsehood" (\textit{De Anima} III.3, 428b18-19).

The definition is redoubled at \textit{De Anima} III.3, 429a1-2, where Aristotle states that "\textit{phantasia} will be movement resulting from an actual exercise of a power of sense." \textit{Phantasia}, then, is \textit{kinēsis} --- not a faculty that stores static images but an ongoing activity in which residual motion is sustained and transformed. This structural feature is decisive. A faculty that merely stored impressions would be a receptacle; \textit{phantasia} is instead a process, an exercise that continues the motion begun in perception into new actualizations that the original perception did not itself produce. $A_0$ established that residual motion is a structural possibility of \textit{kinēsis} --- that motion, once initiated, can persist after its originating cause withdraws. $A_2$ established that the perceptual motion does in fact leave such a residue. \textit{Phantasia} is what does something with that residue.

A crucial distinction must be maintained throughout the analysis: \textit{phantasia} (the activity or faculty) is not identical with the \textit{phantasma} (the image or product). Aristotle preserves the distinction consistently, naming the activity with the abstract noun and the product with the concrete neuter. \textit{Phantasia} is the \textit{kinēsis} "in virtue of which an image is presented to us" (\textit{De Anima} III.3, 428a1-2); the \textit{phantasma} is what appears. O'Gorman captures the relationship precisely: \textit{phantasia} "forms the objects of belief and desire, even gives birth to belief and desire, and provides the images that are the objects of deliberative desire, ethical choice, and political judgment" (\textit{Aristotle's Phantasia}, p. 34). The formation is the work of \textit{phantasia}; the formed image is the \textit{phantasma}.

Collapsing the distinction produces two errors that the subsequent architecture cannot tolerate. The first error treats \textit{phantasia} as a repository of images and thereby loses the \textit{kinēsis} character that Aristotle insists upon. If \textit{phantasia} were a mere storehouse, its characteristic operations --- reactivation of past impressions, recombination of elements, projection of unrealized possibilities, inheritance of affective charge --- would all require some further agency to explain, and \textit{phantasia} would become a passive receptacle rather than an active faculty. The second error treats the \textit{phantasma} as identical with the residual motion on which it depends and thereby loses the ontological distinction that §5 below will develop. If the \textit{phantasma} were merely the residue, it could not be subject to truth and falsity in the way Aristotle insists (\textit{De Anima} III.3, 428b14), for the residue itself is simply the continuation of the original perception and is true or false only in the derivative sense that perception is. The \textit{phantasma}'s capacity for falsity --- its ability to present something that is not there, something that was never there, something that contradicts what the soul otherwise knows to be the case --- requires that it be ontologically distinct from the residue. The distinction is hylomorphic: the residue is material; the \textit{phantasma} is the form that the secondary motion imposes on that material. The residue furnishes the causal substrate; \textit{phantasia} furnishes the activity; the \textit{phantasma} is the actuality in which the two coincide.

\subsection*{3. The Secondary Actualization}

The phantastic motion $M_2$$\to$$A_3$ instantiates the three-factor kinetic schema established at $A_0$, but the factors differ from those at $A_1$$\to$$A_2$. The unmoved originator is no longer the external sensible object; the object has withdrawn, and what remains in its place is the resonant \textit{kinēsis} deposited in the sensory apparatus at $A_2$. This residue --- the dual-trace composite of resonant \textit{aisthēma} and resonant \textit{orexis} --- now functions as the unmoved originator of the phantastic motion: it occasions the secondary actualization without itself being further altered by it. The moved mover is \textit{phantasia} proper, the psychic faculty whose exercise consists in taking up the residual motion and transforming it into the determinate \textit{phantasma}. That which is moved is the central perceptual organ --- the \textit{koinē aisthēsis} or common sense that Aristotle locates at or near the heart --- and, through it, the psycho-somatic state of the perceiver as a whole.

Aristotle's \textit{De Insomniis} provides the physical detail. The residual motions that remain after perception travel inward from the peripheral organs through the sensory apparatus, eventually reaching the central organ where their cognitive significance is realized. The periphery-to-center trajectory is crucial for understanding why the \textit{phantasma} is a genuinely secondary actualization rather than a mere replay of $A_2$. At $A_2$, the actualization occurred at the peripheral organ --- the eye seeing, the ear hearing --- as the joint activity of sensible object and sense faculty. At $A_3$, the actualization occurs at the central organ as the joint activity of the residual motion and \textit{phantasia}. Two distinct locations, two distinct actualizations, two distinct formal results. The first actualization is perception; the second is imagination.

Aristotle marks the difference with a characteristic analogical move. Just as the original sensible object at $A_1$$\to$$A_2$ moved the peripheral organ through the medium, the residual motion at $M_2$$\to$$A_3$ moves the central organ through the blood. Frede captures the continuity and the distinction: the \textit{phantasma} "is the after-image of the residual motion set up in the sense-organ, and once in existence the residual motion has a life of its own." The residue is perception's continuing afterlife in the organ; the \textit{phantasma} is the new actuality that emerges when that afterlife reaches the central organ and there becomes the object of \textit{phantasia}'s activity. Frede observes further that the \textit{phantasma} "must be produced while sense-perception is still in operation to ensure causal continuity; otherwise there would be an unbridgeable causal gap if the \textit{phantasma} or image were not produced while the sense-perception was still in operation" (Frede, \textit{The Cognitive Role of Phantasia}, p. 284). The continuity requirement follows from the three-factor schema: the secondary motion needs an unmoved originator that is genuinely in act, and the residue is genuinely in act only so long as the perceptual motion from which it derives has not fully decayed.

The structural consequence of this architecture deserves emphasis. The phantastic motion is not a psychological add-on but a second instance of the \textit{kinēsis} whose general structure $A_0$ established. As \textit{energeia ateles}, it has its \textit{telos} outside itself --- the actuality it produces is the \textit{phantasma}, which is what building is to the buildable or seeing is to the seen. As relational, it requires both an active potency (the residue as unmoved originator) and a passive potency (the central organ as that which receives). As inheritable, it deposits its own residue --- though the analysis of the \textit{phantasma}'s secondary residues in dreams, memory, and repeated imagination belongs to subsequent sections rather than to the present treatment of how the \textit{phantasma} first emerges. The chain's primitive structure iterates at every stage, and $M_2$$\to$$A_3$ is its second iteration.

\subsection*{4. The Phantasma as Crystallized Actuality}

$A_3$, the \textit{phantasma} proper, is the \textit{energeia} of the secondary motion --- the crystallized image that the central organ presents to the soul as the determinate content of thought, memory, desire, and action. The formulation is decisive. The \textit{phantasma} is not an independent substance that \textit{phantasia} fetches from a mental inventory; it is the completed actuality of the phantastic motion, the terminus at which the secondary \textit{kinēsis} reaches its \textit{telos} and the residual motion becomes a determinate image. Before the central organ is reached, the residue is motion in transit --- real but not yet crystallized. At the central organ, the motion completes, and what had been a mere continuation of perception becomes something new: an image, a determinate representation available for cognitive operations that the residue alone could not sustain.

Three characteristics mark the \textit{phantasma} at its completion. First, it is \textit{stable and recallable}. Where the residue is continuously decaying --- a motion that gradually loses its formal determinacy as the organ returns toward its pre-perceptual state --- the \textit{phantasma} is a completed actuality that persists as a stable image available to the soul across temporal distance from the original perceptual encounter. Aristotle develops the stability at \textit{On Memory} 449b30-450a1, where he observes that memory operates upon images that have been impressed like a seal upon wax, and that the image's persistence depends on the quality of the impression and the condition of the substrate. The stability is not absolute --- images fade, distort, and can be reactivated with varying fidelity --- but it is structurally distinct from the moment-by-moment decay of the residual motion.

Second, the \textit{phantasma} is \textit{formally similar to the original perception yet operates in the absence of the sensible object}. The similarity is guaranteed by \textit{De Anima} III.3, 428b13-14, where Aristotle specifies that the motion of \textit{phantasia} "is necessarily similar in character to the sensation itself." The operation in the object's absence follows from the nature of \textit{phantasia} as a motion produced by actual sensation but capable of persisting beyond it. This dual feature --- similarity to perception combined with independence from the perceived object --- is the source of \textit{phantasia}'s cognitive utility and its liability to error. Because the \textit{phantasma} is similar to what was perceived, it can serve as a cognitive surrogate for the object when the object is absent. Because it can operate in the object's absence, it can misrepresent --- the \textit{phantasma} of the small sun persists even when reason apprehends that the sun is larger than the earth.

Third, the \textit{phantasma} is \textit{the indispensable hinge between perception and thought}. Aristotle states the thesis explicitly at \textit{De Anima} III.7, 431a15-17: "To the thinking soul images serve as if they were contents of perception (and when it asserts or denies them to be good or bad it avoids or pursues them). That is why the soul \textit{never thinks without an image}." The thesis will be developed in §7 below; here it suffices to register its architectural weight. If the \textit{phantasma} is what thought requires in order to operate, then $A_3$ is not a terminal stage but a conditional origin --- the actuality upon which every subsequent cognitive motion in the chain depends.

The \textit{phantasma}'s status as the \textit{energeia} of the secondary motion follows directly from these three characteristics. Stability requires completion; operation-in-absence requires formal independence from the original stimulus; service to thought requires the determinate structure that only a completed actuality possesses. The residue has none of these features in full: it is decaying rather than stable, tied to the organ rather than formally independent, material rather than available as determinate content. The \textit{phantasma} is what emerges when the secondary motion completes these transformations at the central organ.

\subsection*{5. Inheritance of the Dual Trace}

The \textit{phantasma} comes into being already charged. It does not begin as a neutral formal impression that subsequent cognition will colour with affective significance; it is affectively charged from its genesis, because the residue from which it emerges is itself the dual-trace composite established at $A_2$. Resonant \textit{aisthēma} furnishes the formal content that the \textit{phantasma} manifests --- the shape, colour, sound, or scent of what was originally perceived. Resonant \textit{orexis} furnishes the affective valence --- the pursue-or-avoid, the pleasant-or-painful, the good-or-bad as such. Both aspects are inherited; neither is added afterward.

The inheritance is guaranteed by the architecture of §2 and §3 above. Because \textit{phantasia} is a \textit{kinēsis} that arises from the residue and is "necessarily similar in character to the sensation itself" (\textit{De Anima} III.3, 428b13-14), and because the sensation itself --- as §7 of the $A_2$ section argued --- was numerically one and two in being (perception and appetite as co-actualized faculties), the residue preserves both formal aspects and the \textit{phantasma} that emerges from the residue inherits both aspects in turn. The principle operates across every actuality in the chain: what persists preserves what was present at the moment of deposition; what emerges from the persistence inherits what was preserved.

The textual evidence for this inheritance is particularly striking at \textit{Rhetoric} I.11, 1370a27-b1, where Aristotle analyses pleasure in memory and expectation. "Pleasure is the consciousness through the senses of a certain kind of emotion; but imagination is a feeble sort of sensation, and there will always be in the mind of a man who remembers or expects something the imagination of what he remembers or expects. If this is so, it is clear that memory and expectation also, being accompanied by sensation, may be accompanied by pleasure." The pleasure that accompanies memory or expectation is not a fresh evaluation performed on a retrieved formal trace; the image delivers the pleasure directly, because the pleasure was deposited alongside the formal trace in the original perceptual event and has persisted into the present \textit{phantasma}. To remember a former joy \textit{is} to feel a weakened version of that joy; to anticipate a future pain \textit{is} to suffer a weakened version of it in the present. The \textit{phantasma} that memory and expectation deploy is charged with the affective valence the original perception carried, and the charge is delivered as a constitutive feature of the image rather than as an add-on.

Aristotle's account of how \textit{phantasia} moves the body in the absence of external stimuli provides further confirmation. At \textit{De Motu Animalium} 8, 701b17-22 he observes: "This change of quality is caused by imaginations and sensations and by ideas. Sensations are obviously a form of change of quality, and imagination and thinking have the same power as the objects. For in a measure the form conceived, be it of hot or cold or pleasant or fearful, is like what the actual objects would be, and so we shudder and are frightened merely by thinking." The mere image of what is fearful chills the blood; the mere image of what is pleasant heats it. The somatic effects are not generated by an additional evaluative faculty acting upon a neutral image; they are generated by the image itself, because the image carries the affective charge that the original perceptual encounter deposited. The dual-trace inheritance is what makes possible the remarkable phenomenon Aristotle flags at \textit{De Anima} I.1, 403a21-22: "in the absence of any external cause of terror we find ourselves experiencing the feelings of a man in terror." The \textit{phantasma} alone suffices to produce the physiological and affective effects that an external threat would produce, because the \textit{phantasma} is charged with the resonant \textit{orexis} that an external threat would activate.

The inheritance has a structural consequence for the chain's subsequent architecture. If the \textit{phantasma} were a neutral formal trace to which evaluative significance is added downstream, then \textit{doxa} and the higher-order \textit{pathē} would have to generate affective content ex nihilo whenever they engage the image. But the downstream sections do not work this way. \textit{Doxa} ratifies what the \textit{phantasma} presents; it does not add affective valence from outside. The higher-order \textit{pathē} --- fear, anger, pity, shame --- specify and determine the affective valence the \textit{phantasma} already carries; they do not introduce it where it was not before. The dual-trace inheritance at $A_3$ is thus the material condition for the entire cognitive-affective architecture downstream. Without it, the chain could not function.

\subsection*{6. Causal Layering and Ontological Distinctness}

The \textit{phantasma}'s dependence on the residue and its ontological distinctness from the residue require careful specification, because the two features coexist in tension. The \textit{phantasma} is causally downstream of the residue --- it would not exist if \textit{phantasia} had not taken up the residual motion --- yet it is not reducible to the residue, because the \textit{phantasma} possesses features the residue lacks. The clearest analytic framework for this dual relationship comes from Caston's causal-layering analysis of Aristotle's \textit{phantasia}. "Like an echo," Caston observes, "the \textit{phantasma} is only an indirect effect of the object of perception; the \textit{phantasma} is directly caused by the sensory stimulation (\gk{αἴσθημα}), which in turn is directly caused by the object (\gk{αἰσθητόν})" (Caston, \textit{Why Aristotle Needs Imagination}, p. 48). The layering is explicit: the external object causes the \textit{aisthēma}; the \textit{aisthēma} causes the \textit{phantasma}. Three layers, three causal steps, three ontologically distinct items.

The ontological distinctness is evidenced most decisively by the \textit{phantasma}'s capacity for falsity. Aristotle insists that perception of proper sensibles is "always true or admits the least falsehood" (\textit{De Anima} III.3, 428b18-19), but he equally insists that \textit{phantasia} is often false (\textit{De Anima} III.3, 428a11-16). The asymmetry in truth-aptness is the sharpest argument for ontological distinctness: if the \textit{phantasma} were simply identical with the residual trace, it could no more be false than the trace itself, because the trace is simply the continuation of the perceptual motion that recorded the object faithfully. The \textit{phantasma}'s liability to falsity requires that it be something ontologically different --- something to which truth or falsity can attach as a feature it possesses independently of its causal ancestor.

Aristotle's own illustration at \textit{De Anima} III.3, 428b2-4 makes the point vivid. "We imagine the sun to be a foot in diameter though we are convinced that it is larger than the inhabited part of the earth." The \textit{phantasma} of the small sun and the settled belief in the sun's supra-terrestrial magnitude coexist: the image persists unchanged while reason holds a contrary judgment. The coexistence is only possible if the \textit{phantasma} is ontologically distinct from both the judgment that contradicts it and the perceptual residue that generated it. The residue at $A_2$ is what it is: the formal trace of the distal object's apparent size as it impressed upon the eye. The \textit{phantasma} at $A_3$ is a determinate image that can be presented to the soul at a temporal distance from the original encounter and can persist against the soul's own rational correction. Two distinct items, two distinct modes of being.

Kevin White illuminates the ontological specificity of the transition with the observation that "the phantasma points to the transfer of an appearance from its origin in that which appears into another medium" (White, \textit{The Meaning of Phantasia}, p. 487). The "transfer" is precisely what the secondary motion accomplishes: what was an appearance-for-the-organ (the residue lodged in the peripheral sensory apparatus) becomes an appearance-for-the-soul (the \textit{phantasma} presented to the central organ). The medium changes, and with it the mode of being. The residue is a bodily motion in the periphery; the \textit{phantasma} is a determinate formal content available to the soul's cognitive operations. The one is matter-laden; the other is matter-transcending in the sense that it can be recombined, manipulated, and reactivated in ways the residue cannot.

The \textit{phantasma}'s ontological distinctness does not impugn its dependence on the residue. Aristotle is unequivocal: \textit{phantasia} "is impossible without sensation" (\textit{De Anima} III.3, 428b11-12), and the \textit{phantasma} "cannot exist apart from sensation" (\textit{De Anima} III.3, 428b14-15). Without perception, no residue; without residue, no \textit{phantasma}. The dependence is genetic and structural; but it does not entail identity. The \textit{phantasma} is what the residue becomes when it has been taken up by \textit{phantasia} and crystallized at the central organ --- the same affection under a different mode of being, the same \textit{pathos} now available to the soul as an object of further actualization.

One qualification deserves registration. Aristotle allows that \textit{phantasmata} may also be generated from thinking, not only from sense-perception. At \textit{De Motu Animalium} 8, 702a20-21 he writes that "imagination depends either upon thinking or upon sense-perception." The clause registers a dual origin: while sensation is the paradigmatic source of \textit{phantasmata}, thought can also generate images --- as when reason contemplates a geometrical figure that is not presently perceived, or as when practical deliberation projects an unrealized future state. In these cases the \textit{phantasma} is not generated from a fresh residue but is drawn from the mnemonic repertoire and reactivated by thought. The chain's origin at $A_1$ supplies the genetic ground for every \textit{phantasma} ultimately --- no \textit{phantasma} is possible without at some time having been derived from sensation --- but its operative origin in any particular episode may be either perceptual or noetic. The dual origin will become important when the subsequent sections analyse how the cognitive motions at $M_3$$\to$$M_4$ can operate on stored \textit{phantasmata} independently of any ongoing perceptual event.

\subsection*{7. The Phantasma as Object of Thought and Desire}

The architectural role of $A_3$ in the chain is fixed by Aristotle's claim at \textit{De Anima} III.7, 431a15-17: "To the thinking soul images serve as if they were contents of perception (and when it asserts or denies them to be good or bad it avoids or pursues them). That is why the soul never thinks without an image." The claim is redoubled at \textit{De Anima} III.8, 432a3-9: "The soul never thinks without a \textit{phantasma}... the faculty of thought thinks the forms in the \textit{phantasmata}, and as in the former case what is to be pursued or avoided is marked out for it, so where there is no sensation and it is engaged upon the \textit{phantasmata} it is moved to pursuit or avoidance."

Two theses are condensed in these passages. The first --- the cognitive dependence thesis --- is that thought requires \textit{phantasmata} as its material. The soul does not think in a vacuum; it thinks on images, through images, by operating upon images. Even the most abstract intellection proceeds by abstracting the universal from the particular \textit{phantasma}, which means that the \textit{phantasma} is the condition without which intellection could not begin. The second --- the orectic continuity thesis --- is that the soul's pursuit-or-avoidance response to \textit{phantasmata} is structurally continuous with its response to actual sensations. The \textit{phantasma}'s inheritance of resonant \textit{orexis} from the original perceptual event is what makes this continuity intelligible: the \textit{phantasma} can move the soul to pursue or avoid because it carries the affective valence that an actual sensible object would carry, and the soul responds to the charged image as it would respond to the charged object.

The \textit{phantasma}'s role as object of thought and desire grounds the chain's subsequent architecture at $M_3$$\to$$M_4$. The next section will analyse the cognitive motions that operate on $A_3$: intellection (the \textit{phantasma} addressed as bearing a universal), memory (the \textit{phantasma} addressed as likeness of the past), the synthetic mode of deliberative or speculative \textit{phantasia} (multiple \textit{phantasmata} synthesized into a unified object of thought), and the orthogonal committal dimension of \textit{doxa} (the \textit{phantasma}'s aspectual presentation ratified as propositional taking-as-true). Each of these operations presupposes that the \textit{phantasma} is available as a determinate object upon which the rational cognitive apparatus can operate. The present section's architectural contribution is to establish how the \textit{phantasma} comes to be such an object.

The structure is iterative. \textit{Phantasia} at $M_2$$\to$$A_3$ takes up the resonant \textit{kinēsis} deposited at $A_2$ and crystallizes it into the \textit{phantasma} at $A_3$. The cognitive motions at $M_3$$\to$$M_4$ take up the \textit{phantasma} at $A_3$ and produce the completed cognitive actualities at $A_4$ --- the grasped universal, the remembered past, the evaluated possibility, the contemplated truth, or the ratified belief. $A_3$ is the unmoved originator of the cognitive motions in the same way that $A_2$ was the unmoved originator of the phantastic motion. The chain's primitive structure --- actuality grounding motion grounding actuality --- iterates at every stage, and the \textit{phantasma} is the pivotal actuality at which the physical-perceptual register gives way to the cognitive-affective register that the rest of the chain elaborates.

\subsection*{8. Synthesis: $A_3$ as Ground of $M_3$→$M_4$}

$A_3$, then, is the crystallized actuality of the phantastic motion: the \textit{phantasma} proper, arisen from the resonant \textit{kinēsis} deposited at $A_2$ and brought by \textit{phantasia} to completion at the central organ. The section has traced $M_2$$\to$$A_3$ through six moves: the distinction between \textit{phantasia} as \textit{kinēsis} and \textit{phantasma} as product; the application of the three-factor schema to the phantastic motion with its periphery-to-center trajectory; the \textit{phantasma}'s status as crystallized actuality with the features of stability, object-absence, and cognitive hinge; the inheritance of the dual-trace structure such that the \textit{phantasma} comes already charged with both formal content and affective valence; the causal layering that makes the \textit{phantasma} ontologically distinct from the residue while remaining genetically dependent upon it; and the \textit{phantasma}'s architectural role as the object upon which subsequent cognitive and orectic motions operate.

$A_3$'s architectural role in the chain follows the iterated pattern established at $A_0$ and $A_1$$\to$$A_2$. Each completed actuality is simultaneously the terminus of its antecedent motion and the unmoved originator of its successor. $A_3$ is the terminus of the phantastic motion from $A_2$; it is also the unmoved originator of the cognitive motions at $M_3$$\to$$M_4$ through which the completed cognitive actualities at $A_4$ will emerge. The next section takes up those cognitive motions in detail: how \textit{nous} addresses the \textit{phantasma} as bearing a universal; how memory addresses it as a likeness of the past; how the synthetic mode of deliberative and speculative \textit{phantasia} operates on multiple \textit{phantasmata} to construct unified objects of thought that no single \textit{phantasma} contains; and how, orthogonally to all three orientational modes, the committal dimension of \textit{doxa} ratifies the \textit{phantasma}'s aspectual presentation as a propositional truth-claim to which the soul is involuntarily bound.

The reader is accordingly deposited at the threshold of the rational cognitive apparatus --- the faculty whose operations on the \textit{phantasma} generate the completed cognitive actualities from which, in evaluatively complex cases, higher-order \textit{pathē} arise and \textit{orexis} is mobilized into the determinate conative forms that drive appetitive action at $M_3$$\to$$M_4$.\footnote{The Heideggerian analysis of the \textit{phantasma} finds its most direct Aristotelian anchor in the claim that the soul never thinks without an image. Heidegger reads this thesis as a precursor to his own analysis of understanding (\textit{Verstehen}) as always already structured by a pre-predicative "as"-structure that the existential analytic of Dasein makes explicit. The \textit{phantasma}'s inheritance of the dual-trace structure --- and especially its affective charge from resonant \textit{orexis} --- finds a phenomenological parallel in Heidegger's account of \textit{Befindlichkeit} as co-equiprimordial with \textit{Verstehen}: understanding is never neutral reception but is always already structured by the disclosive mood in which Dasein finds itself. A full integration of the Aristotelian \textit{phantasma} with Heidegger's existential account of understanding and state-of-mind lies beyond the scope of the present chapter and is developed in [Chapter X: The Motion-Inducing Soul --- Rhetoric, Attunement, and Actualization (Heidegger, Rickert, Burke)].

\subsection*{Development Note}

Three items flag future development. First, the dual-origin clause at \textit{De Motu Animalium} 8, 702a20-21 --- that \textit{phantasia} depends either upon thinking or upon sense-perception --- raises a structural question for the chain's architecture. If thought can generate \textit{phantasmata} without ongoing perception, then the chain may be re-entered at $A_3$ with stored \textit{phantasmata} reactivated by \textit{nous}, bypassing $A_1$$\to$$A_2$ as operative stages in any particular episode of cognition. The genetic dependency on perception remains absolute --- no \textit{phantasma} is possible without at some time having been derived from sensation --- but the operative origin in any particular episode may be noetic rather than perceptual. The full articulation of this dual-origin structure and its implications for the chain's topology will be developed in the section on cognitive engagement at $M_2$$\to$$M_3$.

Second, the stability and recallability of the \textit{phantasma} established in §4 raises the question of how \textit{phantasmata} persist across extended temporal intervals --- the question, in effect, of memory. Aristotle's \textit{On Memory} develops memory as the activity of attending to a \textit{phantasma} "as a likeness of something past," with the "qua past" structure being what distinguishes memory from the reactivation of a generic \textit{phantasma} without temporal indexing. The relationship between the \textit{phantasma}'s persistence as stable image and its capacity to be addressed as memorial content --- as bearing an indexical relation to a past perceptual encounter --- will be developed in the subsequent section on the orientational modes of cognitive engagement. The relevant primary passages are \textit{On Memory} 449b22-25, 450a11-13, and 450a23, where Aristotle distinguishes memory from imagination and specifies the temporal structure of memorial attention.

Third, the \textit{phantasma}'s ontological distinctness from the residue combined with its genetic dependence on the residue raises a further question about the chain's topology when \textit{phantasmata} are recombined into unified objects of thought --- as in the synthetic mode of deliberative and speculative \textit{phantasia}. When multiple \textit{phantasmata} are synthesized into the image of a state that has never been perceived (the projected future, the geometrical construction, the rhetorically induced picture of the anticipated evil), the resulting composite is ontologically novel but genetically dependent on its constituent \textit{phantasmata}. Aristotle's clearest treatment of this synthesis appears at \textit{De Anima} III.11, 434a5-10, where he specifies that deliberative \textit{phantasia} is "the capacity to make a unity out of several images" --- the faculty that composes novel \textit{phantasmata} from the phantasmatic repertoire. The full analysis of this synthetic capacity, together with its relation to practical reasoning and the sense of time co-implicated with deliberative \textit{phantasia} (\textit{De Anima} III.10, 433b5-10), will be developed in the section on cognitive actualities.}



\end{document}
